\section*{Introduction}

Network analysis is a critical tool for capturing dependencies inherent to the development of social phenomenon. The concept of social network analysis began as early as the 1930s with Jacob Moreno's \textit{Who Will Survive?} and has developed into a rich and diverse field. In Political Science, network analysis has now been utilized to study a diverse range of topics: trade, intergovernmental organizations, sanctions, internal conflict, political behavior.\footnote{For examples, see \cite{dorussen2010}, \cite{cao2009}, \cite{cranmer2014} and for an overview of network analysis in Political Science see \cite{ward2011}.} Yet, unfortunately, a majority of studies in Political Science still rely on the standard dyadic framework. We argue that conflict evolution is a process conditioned on the relative effects of actors' behavior and is thus best conceptualized via a network approach. 	

Sociologists have long established that to understand an actor's behavior you have to understand the context in which they operate as well as interpret their interactions with one partner in light of all interactions across all other partners. Conflict scholars are interested in questions that mirror these exact patterns, such as: which actor is driving the violence within a multi-actor conflict? Did a government crackdown cause anti-government coordination or chaos between armed groups? Are civilian challenges towards non-state armed actors causing an increase or decrease in violence? These questions precisely illuminate the need to apply the appropriate methodological approach to pressing questions in our field.

Our paper, however, is not merely a methodological exercise. We are also motivated by a rich, underdeveloped theoretical question: do nonviolent protests influence the evolution of violence between armed actors?  A recent wave of research in peace studies has illuminated the efficacy of non-violent strategies in contexts of repression. \cite{chenoweth2011} maintains that non-violent campaigns are more likely to succeed in overthrowing a leader than violent campaigns. This research, however, is primary focused on so-called maximalist campaigns wherein an entire population unites under the singular cause of removing a (usually corrupt) leader from power. This previous focus on maximalists movements spurs scholars to explore the capacity of non-violent actions at the local level. 

We are particularly interested in the link between protests and violence in the Mexican case. The Mexican Criminal conflict began in 2006 and has continued to trouble the socio-political landscape of Mexican life. Different estimates place number of drug-related deaths well over 100,000 from 2006-2012. Violence in Mexico is driven by competitive drug trafficking organizations (DTOs) which vie for control for access to resources, transportation routes, and civilian influence. Yet, even in these risky conditions, Mexico has exhibited a rich civic society committed to take action, often through nonviolent protests, to inspire change.

In this paper we utilize an original database of actor-coded conflictual events in Mexico. We investigate whether nonviolent protests influence the likelihood of violence between armed actors. To do so, we generate a conflict network at the quarterly level for the time period 2006-2012 and incorporate novel protest data which documents the number of protests against violence at the state level during the Mexican Criminal Conflict. 